\section*{الفصل \ref{ch:b3:field-lab-methods-3} --- طرق الحقل والمختبر III}
% Values from the Book 5 scratch calculator calc/ch31.py.

\begin{solution}{exo:b3:field-lab-methods-3:1}
(أ)~$\qty{0.01}{\milli\gal} = \qty{e-7}{m.s^{-2}}$،
و$\num{e-7}/\num{9.81} = \num{1.0e-8}$: جزء واحد من مئة مليون.

(ب)~$0.01/0.3086 = \qty{0.032}{m}$، أي نحو \qty{3.2}{cm}.

(ج)~يغيّر التأثيران معًا الجاذبية بمقدار
\qty{0.1967}{\milli\gal.m^{-1}}، وهو الفرق $0.3086 - 0.1119$، لذا
يجب أن يتغيّر الارتفاع بمقدار $0.01/0.1967 = \qty{5.1}{cm}$.
\end{solution}

\begin{solution}{exo:b3:field-lab-methods-3:2}
$\Delta m = \tfrac43\pi(\qty{50}{m})^3\times\qty{500}{kg.m^{-3}} =
\qty{2.62e8}{kg}$. تبلغ القمّة $G\Delta m/z^2 =
\num{6.674e-11}\times\num{2.62e8}/\num{e4} = \qty{1.75e-6}{m.s^{-2}}$,
أي \qty{0.175}{\milli\gal}. نصف العرض هو $x_{1/2} = z/1.305 =
\qty{77}{m}$.
\end{solution}

\begin{solution}{exo:b3:field-lab-methods-3:3}
(أ)~بفرض $AM = MN = NB = a$، فإنّ $K = 2\pi/(1/a - 1/2a - 1/2a + 1/a) =
2\pi a = \qty{62.83}{m}$، و$\rho_a = 62.83\times0.239/0.1 =
\qty{150}{\ohm\metre}$. (ب)~$K = \pi(15^2 - 5^2)/(2\times5) = \pi\times
200/10 = \qty{62.83}{m}$. تقع الأقطاب عندئذٍ عند $0$، $10$، $20$
و\qty{30}{m}، بالترتيب AMNB: وهي ترتيب فينر ذو تباعد
\qty{10}{m}.
\end{solution}

\begin{solution}{exo:b3:field-lab-methods-3:4}
(أ)~$V = 0.2998/\sqrt5 = \qty{0.134}{m.ns^{-1}}$. (ب)~$t =
2\times10/0.134 = \qty{149}{ns}$. (ج)~$\delta = 0.0053\sqrt5/\num{e-3} =
\qty{11.9}{m}$، وعمق مفيد نحو \qty{36}{m}. (د)~في الطين
$\delta = 0.0053\sqrt{20}/0.01 = \qty{2.4}{m}$، وعمق مفيد نحو
\qty{7}{m}، وسرعة لا تتعدّى \qty{0.067}{m.ns^{-1}}. ويبقى
التقريب صالحًا هنا: عند \qty{100}{MHz}،
$\omega\varepsilon = 2\pi\times\num{e8}\times20\times\num{8.85e-12} =
\qty{0.11}{S.m^{-1}}$، أي عشرة أمثال $\sigma$.
\end{solution}

\begin{solution}{exo:b3:field-lab-methods-3:5}
(أ)~تملك الأسطوانة $\lambda = \pi R^2\Delta\rho$ لكلّ وحدة طول،
ووخارجها تكون جاذبيتها كجاذبية كتلة خطية على محورها (مبرهنة غاوس
على أسطوانة متحدة المحور): $\Delta g = 2G\pi R^2\Delta\rho\,z/(x^2
+ z^2)$. يبلغ نصف القمّة عندما $x^2 + z^2 = 2z^2$، إذن $z =
x_{1/2}$. (ب)~تقع الكرة عند $1.305\times100 = \qty{130}{m}$،
والأسطوانة عند \qty{100}{m}: المقطع نفسه، إذا قُرئ ككرة، يضع الجسم
أعمق بنسبة \qty{30}{\percent}. يتناقص شذوذ الكرة أسرع مع
المسافة، لذا يجب أن تكون الكرة أعمق لتكون بالاتساع نفسه.
\end{solution}

\begin{solution}{exo:b3:field-lab-methods-3:6}
(أ)~تمثّل كل نقطة مساحة $\qty{100}{m}\times\qty{100}{m} =
\qty{e4}{m^2}$، إذن $\iint\Delta g\,\mathrm{d}A \approx
\num{-438e-5}\times\num{e4} = \qty{-43.8}{m^3.s^{-2}}$، و$\Delta m
= -43.8/(2\pi\times\num{6.674e-11}) = \qty{-1.04e11}{kg}$. (ب)~$1.04/1.74
= \qty{60}{\percent}$. يقع الباقي في أطراف الشذوذ خارج
الخريطة: يصل الغرانيت إلى عمق يزيد على كيلومتر،
وشذوذ جزئه العميق يتناقص ببطء مع المسافة. ويؤدي توسيع الخريطة
بما يكفي، أو بنموذج، إلى استعادته.
\end{solution}

\begin{solution}{exo:b3:field-lab-methods-3:7}
(أ)~بالنسبة إلى القاطع الرفيع، يهبط الشذوذ إلى نصف قمّته عند $x = \pm z$:
يبلغ عرضه الكامل عند نصف القمّة $2z = \qty{60}{m}$،
وثلاث قراءات عبره تحتاج إلى تباعد لا يزيد على \qty{20}{m}. (ب)~$2\times
220/3 = \qty{147}{m}$. (ج)~يتناسب تباعد المحطات المطلوب مع
عمق أقرب جسم مهم: إن مسحًا مصمَّمًا للقاطع العميق
عند \qty{150}{m}، سيُشوّه شذوذ القاطع الضحل، الذي
لا يزيد عرضه على \qty{60}{m}؛ والمسح المصمَّم للجسم الضحل يرى
القاطع العميق أيضًا، بقراءات كثيرة تفوق ما يحتاجه.
\end{solution}

\begin{solution}{exo:b3:field-lab-methods-3:8}
(أ)~$F = 80/8 = 10$ و$\phi = F^{-1/m} = 10^{-1/2} = \num{0.316}$.
(ب)~$10^{-1/1.8} = \num{0.278}$ و$10^{-1/2.2} = \num{0.351}$. ويؤدي تغيّرٌ
قدره \num{0.2} في $m$ إلى تحريك المسامية بنحو أربع نقاط مئوية؛
أما قيم المقاومية، المقاسة بدقة بضعة بالمئة، فتحرّكها أقل بكثير. أُسّ
التلاحم يتحكّم في عدم اليقين، ويجب معايرته
على عيّنات من الصخر نفسه.
\end{solution}

\begin{solution}{exo:b3:field-lab-methods-3:9}
(أ)~عند $s = 2h_1$، و$(2nh_1/s)^2 = n^2$، فتكون الحدود $k^n(1 +
n^2)^{-3/2}$: $0.3536k$، $0.0894k^2$، $0.0316k^3$. عند $k = 0.5$:
$\num{0.1768}$، $\num{0.0224}$، $\num{0.0040}$، و$\rho_a/\rho_1 \approx
1 + 2\times0.2031 = \num{1.406}$ (يعطي المجموع الكامل \num{1.409}). عند $k
= -0.5$ تتناوب الحدود في الإشارة، $-0.1768 + 0.0224 - 0.0040$،
و$\rho_a/\rho_1 \approx \num{0.683}$ (المجموع الكامل \num{0.685}). (ب)~عند $s \ll
h_1$ تتلاشى كل الحدود و$\rho_a \to \rho_1$، وعند $s \gg h_1$ يؤول كل
قوس إلى~1 و$1 + 2\sum k^n = 1 + 2k/(1 - k) = (1 + k)/(1 -
k) = \rho_2/\rho_1$. (ج)~$\rho_2/\rho_1 = (1 + k)/(1 - k) = 3$ و$1/3$.
\end{solution}

\begin{solution}{exo:b3:field-lab-methods-3:10}
(أ)~بالترتيب AMNB وبتباعد $a$: $AM = a$، $BM = 2a$، $AN = 2a$،
$BN = a$، فيكون القوس $2/a - 1/a = 1/a$ و$K = 2\pi a$.
(ب)~بوضع A وB عند $\mp s$ وM وN عند $\mp l$: $AM = BN = s - l$،
$BM = AN = s + l$، ويكون القوس $2/(s - l) - 2/(s + l) = 4l/(s^2 -
l^2)$، و$K = \pi(s^2 - l^2)/(2l)$. عندما $l \ll s$، فإنّ $\Delta V_{MN}
\approx 2l\,E$ حيث $E = 2(-\mathrm{d}V/\mathrm{d}r)$ عند المركز، أي
حقل القطبين الكهربائيين؛ ثم $\rho_a = K\Delta V/I \to
\pi s^2/(2l)\times 4l(-\mathrm{d}V/\mathrm{d}r)/I = 2\pi s^2(-
\mathrm{d}V/\mathrm{d}r)/I$. (ج)~نسبة العامل التقريبي إلى
العامل الدقيق هي $s^2/(s^2 - l^2) = 1/(1 - 0.01) = 1.010$: أي واحد بالمئة.
\end{solution}

\begin{solution}{exo:b3:field-lab-methods-3:11}
(أ)~القراءة: \qty{0.010}{\milli\gal}. الارتفاع:
$0.1967\times0.05 = \qty{0.010}{\milli\gal}$. الكثافة: تصحيح بوغير
هو $0.04191\,h\,d$ مع $d$ بالوحدة \unit{g.cm^{-3}}، لذا
$0.04191\times50\times0.1 = \qty{0.210}{\milli\gal}$. التضاريس:
\qty{0.200}{\milli\gal}. المجموع: $\sqrt{0.010^2 + 0.010^2 + 0.210^2 +
0.200^2} = \qty{0.29}{\milli\gal}$. (ب)~تهيمن الكثافة والتضاريس،
كلٌّ منهما نحو عشرين ضعف خطأ القراءة. ولرصد
\qty{0.1}{\milli\gal}، يجب قياس كثافة اللوح المندس على
عيّنات أو تُضبط بملاءمتها للبيانات (القيمة التي تعكس الطبوغرافيا بأقل قدر
ممكن)، ويُصحَّح تأثير التضاريس انطلاقًا من نموذج ارتفاع دقيق.
(ج)~$\sqrt{0.010^2 + 0.010^2 + 0.210^2} = \qty{0.21}{\milli\gal}$.
\end{solution}

\begin{solution}{exo:b3:field-lab-methods-3:12}
(أ)~الكتلة التي تطابق القراءة المركزية هي $\Delta m = \Delta g(0)z^2/G
= \num{-4.16e-5}z^2/\num{6.674e-11}$: \qty{-5.61e10}{kg} عند
\qty{300}{m}، \qty{-9.97e10}{kg} عند \qty{400}{m}، \qty{-1.56e11}{kg} عند
\qty{500}{m}. القراءات المتوقعة عند $x = \pm300$ و$\pm600$ هي
\num{-1.47} و\num{-0.37}، \num{-2.13} و\num{-0.71}، \num{-2.62}
و\qty{-1.09}{\milli\gal}. وقيم \omterm{def:b3:field-lab-methods-3:inversion-practice}{عدم التطابق} هي $\chi^2 = 1238$، $3.7$
و$856$. (ب)~لا يقع ضمن حدود الخطأ إلا النموذج عند \qty{400}{m}، مع $\chi^2 = 3.7$ لخمس
بيانات. (ج)~نعم: أيّ جسم متناظر كرويًا
بالكتلة نفسها والمركز نفسه يعطي الشذوذ نفسه تمامًا، ويمكن لأشكال أخرى أن تقع
ضمن حدود الأخطاء. كثافة تُقاس على عيّنات، أو في ثقب حفر، أو بطريقة أخرى
هي التي ستحسم الأمر.
\end{solution}

\begin{solution}{pb:b3:field-lab-methods-3:1}
\textbf{1.} $(2846.708 - 2846.512)/\qty{2.0}{h} = \qty{0.098}{\milli\gal.h^{-1}}$.

\textbf{2.} $2851.075 - 0.098\times0.5 - 2846.512 = \qty{4.514}{\milli\gal}$;
$2850.497 - 0.088 - 2846.512 = \qty{3.897}{\milli\gal}$;
$2849.392 - 0.127 - 2846.512 = \qty{2.753}{\milli\gal}$.

\textbf{3.} الارتفاعات فوق القاعدة هي $-27.1$، $-30.6$
و\qty{-18.1}{m}؛ وبإضافة $0.1967\,\Delta h$ تُعطي شذوذات بوغير
$\num{-0.817}$، $\num{-2.122}$ و\qty{-0.808}{\milli\gal}. المحطات
أدنى من القاعدة وأقرب إلى مركز الأرض،
لذا تكون قراءاتها الخام أعلى؛ ويزيل التصحيح ذلك
ويُبقي انخفاضًا فوق الغرانيت.

\textbf{4.} الفرق الإضافي $0.098 - 0.040 = \qty{0.058}{\milli\gal.h^{-1}}$،
\qty{0.12}{\milli\gal} على امتداد الحلقة هو المدّ الأرضي الصلب. خلال
ساعتين يتغيّر تقريبًا تغيّرًا خطيًا، لذا يُزال مع الانجراف بمستقيم يمرّ بقراءتي
القاعدتين.

\textbf{5.} الأخطاء هي $\num{0.008}$، $\num{0.009}$
و\qty{0.017}{\milli\gal}. قد يجعل تقريب كل ارتفاع إلى \qty{0.1}{m} فرقَ الارتفاع
خاطئًا بمقدار يصل إلى \qty{0.1}{m}، أي
\qty{0.020}{\milli\gal}: يفسّر التقريب الأخطاء، أمّا المدّ
المتبقي بعد المستقيم فدون \qty{0.001}{\milli\gal}.

\textbf{6.} $0.01/0.1967 = \qty{5.1}{cm}$. أمّا \qty{3.3}{cm} الواردة في الكتيّبات فتساوي
$0.01/0.3086$: لا يتعلّق شذوذ الهواء الحر بالكثافة
المختارة لتصحيح بوغير، ويجب أن يخدم الارتفاع ذلك أيضًا.

\textbf{7.} $x_{1/2} = (3.952 - 3.048)/2 = \qty{0.452}{km}$، و$z = 1.305
\times452 = \qty{590}{m}$ تحت سطح الأرض.

\textbf{8.} $\Delta m = \Delta g\,z^2/G = \num{-2.293e-5}\times590^2/
\num{6.674e-11} = \qty{-1.20e11}{kg}$. مع $\Delta\rho =
\qty{-130}{kg.m^{-3}}$ يبلغ حجمه \qty{9.2e8}{m^3} ونصف قطره
\qty{603}{m}، أكبر من عمق مركزه، فالكرة ستبرز
من سطح الأرض. لا توجد كرة بتباين كثافة الغرانيت تطابق
القياسات؛ فيجب أن يبلغ الجسم قرب السطح.

\textbf{9.} يمتدّ الغرانيت من السطح إلى عمق كبير. أمّا جزؤه الضحل
الواقع قرب المحطات فيجعل القمّة ضيّقة، والقاعدة التي
لا تعرف إلا العرض، تضع كرةً واحدة عند عمق متوسط؛ والجزء العميق
يضيف جوانب عريضة تفتقر إليها الكرة (انظر شكل المقطع)، والكتلة التي تحويها
تغيب عن الحساب.

\textbf{10.} $2\pi G\Delta\rho = 2\pi\times\num{6.674e-11}\times(-130) =
\qty{-5.45e-8}{s^{-2}}$. مع $z_1 = 0$، $z_2 = \qty{1645}{m}$، $a =
\qty{461}{m}$: القوس يساوي $1645 + 461 - 1708.4 = \qty{397.6}{m}$،
و$\Delta g = \qty{-2.17}{\milli\gal}$. أمّا عند عمق غير محدود فيؤول القوس
إلى $a$، و$\Delta g = \qty{-2.51}{\milli\gal}$.

\textbf{11.} يجب أن يكون القوس $2.29/\num{5.45e-3} = \qty{420}{m}$، وهو ما يعطيه
قاع عند نحو \qty{2.6}{km}. وأيّ قاع بين نحو
\qty{1.6}{km} وعمقٍ غير محدود يغيّر القمّة بأعشار قليلة
من الميليغال فقط، وجسم أعرض قليلًا أو أكثف يعطي القمّة نفسها
بقاع أقرب إلى السطح، فالقاع غير محسوم.

\textbf{12.} مع $z_1 = \qty{6}{m}$: $1639 + 461.0 - 1708.4 = 391.6$،
و$\Delta g = \qty{-2.14}{\milli\gal}$، بأقلّ بمقدار \qty{0.03}{\milli\gal} من
حالة القمّة عند سطح الأرض. وهو دون أخطاء التضاريس والكثافة
في المسح: لا تستطيع الجاذبية أن تحدّد هل يصل الغرانيت
إلى السطح تحت الطمي. أمّا السبر فيستطيع ذلك.

\textbf{13.} $K = \pi(20^2 - 1^2)/(2\times1) = \qty{626.7}{m}$، و$\Delta V =
\rho_a I/K = 821\times0.050/626.7 = \qty{65.5}{mV}$.

\textbf{14.} $V = 2\times1/\sqrt{17.1^2 - 10.7^2} = 2/13.34 =
\qty{0.150}{m.ns^{-1}}$، $\varepsilon_r = (0.2998/0.150)^2 = 4.0$، أي
رمل جاف؛ والأنبوب على عمق $0.150\times10.7/2 = \qty{0.80}{m}$.

\textbf{15.} $0.150\times20.0/2 = \qty{1.50}{m}$.

\textbf{16.} $20/821 = \num{0.0244}$, $30/1309 = \num{0.0229}$, $40/1684 =
\num{0.0238}$: $S \approx \qty{0.0237}{S}$.

\textbf{17.} تُسهم الطبقة الجافة بمقدار $1.5/600 = \qty{0.0025}{S}$، أمّا
المشبعة فبمقدار $\qty{0.0212}{S}$، لذا يكون سمكها $0.0212\times222 =
\qty{4.7}{m}$، ويقع الغرانيت على عمق $1.5 + 4.7 = \qty{6.2}{m}$ تحت
القاع. التقدير مرتفع: الغرانيت ليس لا نهائي المقاومية، لذا عند هذه المسافات
ما زال $\rho_a$ دون الخط
$s/S$، و$s/\rho_a$ يُبالغ في تقدير $S$. (يتضمّن نموذج المنطقة
\qty{6}{m} من الطمي.)

\textbf{18.} $\phi = (20/222)^{1/2} = \num{0.30}$. مع $m = 1.3$، $\phi =
0.0901^{1/1.3} = \num{0.16}$؛ ومع $m = 2.8$، $\phi = \num{0.42}$.

\textbf{19.} بطرح القاعدة من الخط نحصل، من الغرب إلى الشرق:
3.4, 4.1, 4.8, 5.9, 7.4, 9.2, 11.5, 14.2, 17.2, 19.8, 21.0, 20.1, 17.6,
14.5, 11.5, 9.1, 7.2, 5.7, 4.6, 3.8 و\qty{3.1}{nT}.

\textbf{20.} قبل التصحيح، تبلغ القمّة \qty{28.0}{nT} عند \qty{2.20}{km}،
ونصفها، \qty{14.0}{nT}، لا يُبلغ أبدًا إلى الشرق: ينتهي الخط
عند \qty{15.2}{nT}، لأنّ الحقل ارتفع بمقدار \qty{12}{nT} أثناء
السير.

\textbf{21.} القمّة هي \qty{21.0}{nT} عند \qty{2.20}{km}؛ ونصفها هو
\qty{10.5}{nT}. غربًا: بين \num{9.2} عند \num{1.95} و\num{11.5} عند
\qty{2.00}{km}، عند $1.95 + 0.05\times1.3/2.3 = \qty{1.978}{km}$. شرقًا:
بين \num{11.5} عند \num{2.40} و\num{9.1} عند \qty{2.45}{km}، عند
$2.40 + 0.05\times1.0/2.4 = \qty{2.421}{km}$. نصف العرض هو $(2.421
- 1.978)/2 = \qty{0.221}{km}$.

\textbf{22.} الشذوذ المتناظر في حقل مائل بزاوية \ang{65} يعني أنّ
المغنطة الرأسية وحدها هي المؤثّرة: يمتدّ الجسم على امتداد خط الزوال
المغناطيسي، شمال--جنوب. وتفترض القاعدة الواردة في
\cref{prop:b3:field-lab-methods-3:dyke} أنّ القاطع رقيق، عمودي،
ثنائي البعد يمتدّ عميقًا إلى الأسفل، ممغنط بالحثّ، وتعمل قمّته
بوصفها خطًّا من الشحنة المغناطيسية.

\textbf{23.} $\kappa t = 2\pi z\,\Delta T_\mathrm{max}/(F\sin^2 I) = 2\pi\times
221\times21.0/(48000\times0.8214) = \qty{0.74}{m}$. مع $\kappa =
0.03$، $t = \qty{25}{m}$. الجداء وحده هو المحدّد، فقاطع ضعف
السماكة بنصف القابلية المغناطيسية يعطي الشذوذ نفسه.

\textbf{24.} عرض الشذوذ عند نصف قمّته \qty{443}{m}، والمحطات
تبعد \qty{50}{m} عن بعضها: تسع قراءات عبره، أي ثلاثة أمثال
الحدّ الأدنى الذي تشترطه القاعدة.

\textbf{25.} قمّة القاطع
\textbf{\qty[text-series-to-math=true]{220}{m}} تحت سطح الأرض. ويضع النموذج
الجيولوجي عدم التوافق على عمق \qty{222}{m} عند هذه النقطة:
ويرتفع القاطع حتى قاعدة الغطاء ويُقطع عندها، وهو أقدم
من الحجر الرملي الترياسي الذي يدفنه.
\end{solution}
